\subsection{Conteo de Inversiones con Fenwick Tree}

\graybold{Preguntas guía:}

\begin{itemize}
\item ¿Necesito contar, para una secuencia de pares (o permutación), cuántas parejas están "invertidas" (fuera de orden relativo)?
\item ¿El problema se puede reformular como comparar el orden de un mismo conjunto visto de dos formas distintas (por ejemplo, en dos extremos de una figura)?
\item ¿N es grande (\textasciitilde{}10\textsuperscript{5}), descartando el enfoque \bigO{n^2} de comparar todas las parejas directamente?
\end{itemize}

\graybold{¿Para qué sirve?} Contar en \bigO{n log n} el número de pares fuera de orden (inversiones) en una secuencia, una subrutina clásica que aparece en problemas de conteo combinatorio, distancia de Kendall-tau, y — como en este ejercicio — en fórmulas de conteo de regiones en arreglos geométricos.

\graybold{Complejidad:} \bigO{n log n} (ordenar + un Fenwick Tree de tamaño n).

\graybold{Fundamento teórico:} Se recorre la secuencia insertando elementos en un Fenwick Tree indexado por su VALOR (comprimido a rango $[1,n]$); al insertar el elemento actual, la cantidad de elementos YA insertados con valor MAYOR (calculada restando la consulta de prefijo del Fenwick al total ya insertado) es exactamente la cantidad de inversiones que ese elemento forma con los anteriores. Sumando esto en cada inserción se obtiene el total en \bigO{n log n}. Para el problema de la pizza, el número de regiones resultantes de H cortes horizontales y V verticales que pueden cruzarse entre sí (incluso entre cortes del mismo tipo) es:
$$1 + H + V + H \cdot V + \text{inv}(H) + \text{inv}(V)$$
donde cada corte cruza a TODOS los del otro tipo exactamente una vez (garantizado geométricamente, ya que uno va de borde a borde en una dirección y el otro en la perpendicular), y $\text{inv}(\cdot)$ cuenta cuántas parejas de cortes del MISMO tipo invierten su orden relativo entre el lado izquierdo/inferior y el derecho/superior de la pizza — justamente porque dos curvas simples que van de un lado al opuesto se cruzan si y solo si su orden relativo se invierte.

\graybold{Ejercicio aplicado}

Una pizza rectangular recibe H cortes de izquierda a derecha y V cortes de abajo hacia arriba (cada uno puede cruzar a cualquier otro, incluidos los de su mismo tipo). Dadas las coordenadas de entrada/salida de cada corte, calcular en cuántas piezas queda dividida la pizza.

\graybold{I/O:} H,V ≤ 10\textsuperscript{5} → lectura total con tokenización (patrón 2).

\graybold{Código solución}

\begin{lstlisting}[language=Python]
import sys

def count_inversions(pairs):
    # ordena por el primer elemento del par, y cuenta inversiones en el
    # segundo elemento (cuantas parejas quedan "fuera de orden")
    pairs.sort()
    rights = [p[1] for p in pairs]
    sorted_r = sorted(set(rights))
    comp = {v: i + 1 for i, v in enumerate(sorted_r)}   # compresion de coordenadas
    n = len(sorted_r)
    tree = [0] * (n + 1)

    def update(i):
        while i <= n:
            tree[i] += 1
            i += i & (-i)

    def query(i):
        s = 0
        while i > 0:
            s += tree[i]
            i -= i & (-i)
        return s

    inv = 0
    cnt = 0
    for r in rights:
        c = comp[r]
        inv += cnt - query(c)     # ya insertados con 'right' MAYOR = inversion con este
        update(c)
        cnt += 1
    return inv

def solve():
    data = sys.stdin.buffer.read().split()
    idx = 0
    X = int(data[idx]); Y = int(data[idx + 1]); idx += 2
    H = int(data[idx]); V = int(data[idx + 1]); idx += 2
    hcuts = []
    for _ in range(H):
        y1 = int(data[idx]); y2 = int(data[idx + 1]); idx += 2
        hcuts.append((y1, y2))
    vcuts = []
    for _ in range(V):
        x1 = int(data[idx]); x2 = int(data[idx + 1]); idx += 2
        vcuts.append((x1, x2))

    inv_h = count_inversions(hcuts)
    inv_v = count_inversions(vcuts)
    # +1 (pizza entera) +H+V (cada corte agrega al menos 1 region) +H*V
    # (cada corte H cruza a cada corte V exactamente una vez) + inversiones
    # (cruces adicionales entre cortes del MISMO tipo)
    regions = 1 + H + V + H * V + inv_h + inv_v
    print(regions)

solve()
\end{lstlisting}
